\documentclass[11pt]{article}
\usepackage{preamble}
\renewcommand{\answer}[1]{}

\begin{document}

\ladderheading{AP Statistics \textperiodcentered\ Unit 4 \textperiodcentered\ Topic 4.3 \textperiodcentered\ B: 2027-01-15 \textperiodcentered\ E: 2027-03-03}{One Interval, Two Cutoffs}

\focusbanner{TODAY'S PROBLEM}

Suppose a library wants the mean setup time for all 800 kiosk transactions in a month below 45 seconds and hopes to advertise a mean below 42 seconds. An SRS of 25 has the shown summary. A 95\% one-sample $t$-interval is $(39.923,44.877)$ seconds.\par\smallskip \textbf{Mina:} ``I would compare each cutoff with the full interval and keep the conclusion about the population mean.''\par \textbf{Gabe:} ``Because the interval contains times below both cutoffs, I think it supports both claims; I also read 95\% as the share of transaction times inside it.''\par
\begin{center}
\textbf{Kiosk sample}\par\smallskip
\begin{tabular}{|l|c|}
\hline
\textbf{Summary} & \textbf{Value} \\ \hline
\textbf{Population} & 800 \\ \hline
\textbf{$n$} & 25 \\ \hline
\textbf{Mean (s)} & 42.4 \\ \hline
\textbf{SD (s)} & 6.0 \\ \hline
\textbf{Shape} & Symmetric; no outliers \\ \hline
\end{tabular}
\end{center}
\begin{firsttakebox}{FIRST TAKE --- before the video (5 min)}
Does the range 39.923 to 44.877 make both time claims believable? Give one reason.
\workspace{0.9in}
\end{firsttakebox}

\begin{callout}{\IconBook\ LET THE ENTIRE INTERVAL GOVERN THE CLAIM (keep on the page)}{calloutgreen}
A $C\%$ confidence interval estimates a population mean, not the spread of individual observations.
We are $C\%$ confident that the interval captures the population mean represented by the random sample;
the fixed interval either contains that mean or it does not. A value inside the interval remains plausible.
To support $\mu<c$, the entire interval must lie below $c$; an endpoint merely falling below $c$ is not enough.
With other features held fixed, a larger sample tends to narrow a mean interval, while a higher confidence
level widens it.
\end{callout}

\newpage
\textbf{Finish after the video:}\par\smallskip

\dokbadge{1}~\textbf{(a)} Locate 42 and 45 relative to the interval.\par
\workspace{0.7in}

\dokbadge{2}~\textbf{(b)} Interpret the supplied 95\% interval in one sentence about the mean setup time for all 800 transactions.\par
\workspace{1.4in}

\dokbadge{3}~\textbf{(c)} In 3--4 sentences, decide whose approach is justified. Use the interval to judge each cutoff, then explain what Gabe gets wrong about individual transaction times.\par
\workspace{2.0in}

\begin{sentenceframebox}\raggedright\textbf{Frame:} Relative to the interval, 42 is \blank[1.0in] and 45 is \blank[1.0in]; together, those locations support \blank[2.5in].\end{sentenceframebox}

\begin{sentenceframebox}\raggedright\textbf{Frame:} The interval estimates \blank[2.0in] rather than \blank[1.9in], so the advertisement should \blank[2.0in].\end{sentenceframebox}

\begin{summaryexitbox}{EXIT --- turn this in}
One thing the video changed about my first take: \hrulefill\par
\workspace{0.4in}
\end{summaryexitbox}

\end{document}
